% 02_methods.tex - Seccion de metodos (PRISMA 2020, items 5-15).
\section{Methods}

\subsection{Protocol and registration}
% TODO: protocolo registrado ANTES de ejecutar la busqueda (PROSPERO u OSF).
% La URL se registra en data/stats_summary.json (study_metadata.preregistration_url).
\todo{DATO PENDIENTE: preregistration_url de data/stats_summary.json}

\subsection{Eligibility criteria}
% TODO: criterios de inclusion y exclusion definidos en el protocolo.

\subsection{Information sources}
% TODO: bases de datos consultadas y justificacion de cada inclusion o
% exclusion (PubMed, Scopus, IEEE Xplore, ...). La estrategia completa se
% documenta en rag/calls/ ANTES de ejecutar la busqueda (prisma-synthesis F1).

\subsection{Search strategy}
% TODO: cadenas booleanas, campos consultados (titulo/resumen/keywords),
% filtros y rango de fechas.

\subsection{Selection process}
% TODO: deduplicacion, cribado por titulo y resumen, cribado a texto completo,
% numero de revisores independientes y resolucion de desacuerdos.
% Los recuentos de cada etapa se reportan en src/figures/prisma_flow.tex.

\subsection{Data collection process}
% TODO: extraccion con la hoja data/extraction_sheet.csv (columnas
% documentadas en context/prisma_checklist.md). Cada extracto se registra
% ademas en context/evidence_notes.md con su citekey.

\subsection{Risk of bias assessment}
% TODO: herramienta por diseno (RoB 2 para ensayos aleatorizados, ROBINS-I
% para observacionales) y veredicto por estudio (prisma-synthesis F4).

\subsection{Synthesis methods}
% TODO: sintesis narrativa estructurada o meta-analisis; si es cuantitativa,
% medidas de efecto, modelo de efectos fijos o aleatorios y evaluacion de la
% certeza por resultado con GRADE (prisma-synthesis F5).
